\chapter{Родитель паринга Вайнберга и граница нейтринной амплитуды}
\label{ch:v6-spectral-transition-weinberg-pairing-parent-gate}

\section{Цель ретроспективного замыкания}

Предыдущий гейт заменил сингулярную нормированную нейтринную линию
регулярными квадратичными ковариантами
\begin{equation}
 W_\nu(H)=\widetilde H\widetilde H^\dagger,
 \qquad
 B_\nu(H)=\widetilde H\widetilde H^T.
\end{equation}
Они фиксируют слабое направление, но ещё не задают фермионный паринг.
Теперь проверяется сильное утверждение: выводит ли уже существующий
родитель полный оператор Вайнберга, включая семейный тензор,
коэффициент и подавляющий масштаб?

Нулевая гипотеза утверждает, что gauge-ковариантность, голономия $C_3$,
проектор $P_0$ и вес $1/7$ однозначно определяют паринг без нового
спектрального момента. Стоп-критерием служит непрерывное семейство
допустимых мер с одинаковой кинетикой и различным парингом.

\section{Операторный тип существует}

На $H_{15}$ нет правого нейтрино. Минимальный gauge-инвариантный
нейтринный массовый оператор имеет вид
\begin{equation}
 \mathcal O_5=
 \frac{C_{ij}}{\Lambda}
 (\overline{L_i^c}\,\widetilde H^*)
 (\widetilde H^\dagger L_j),
 \qquad C^T=C.
 \label{eq:v6-weinberg-pairing}
\end{equation}
Он имеет размерность пять и меняет лептонное число на две единицы
\cite{weinberg-parent-weinberg1979}. Слабое свёртывание одномерно и
совпадает с ковариантом $B_\nu(H)$. Следовательно, тип оператора
допустим, но допустимость не равна происхождению коэффициента.

Фермионное спектральное действие действительно способно породить
хиггс-квадратичный майорановский член в почти коммутативной геометрии без
правого нейтрино \cite{weinberg-parent-sakellariadou-sitarz2019}.
Однако этот механизм использует нескалярную чётную часть спектральной
функции, которая является дополнительным данными меры.

\section{Семейная классификация}

Для трёхциклической голономии
\begin{equation}
 C_3=
 \begin{pmatrix}0&0&1\\1&0&0\\0&1&0\end{pmatrix}
\end{equation}
условия майорановской симметрии дают
\begin{equation}
 C^T=C,
 \qquad C_3^TCC_3=C.
\end{equation}
Пространство решений двумерно:
\begin{equation}
 C=xP_0+y(I-P_0),
 \qquad
 P_0=\frac13(I+C_3+C_3^2).
 \label{eq:v6-weinberg-family-tensor}
\end{equation}
Машинный аудит воспроизводит размерность два. Поэтому одна $C_3$-
симметрия не фиксирует полный семейный тензор.

Сжатие на уже выведенную нулевую ветвь даёт
\begin{equation}
 P_0CP_0=xP_0.
\end{equation}
Размерность направления становится равной единице с остатком
$1.67\cdot10^{-16}$. Это фиксирует семейную ось одного низкоэнергетического
канала, но не комплексную амплитуду $x$ и не полный трёхсемейный спектр.

\section{Явный no-go спектральной меры}

Рассмотрим допустимое семейство быстро убывающих функций
\begin{equation}
 g_\alpha(z)=ze^{-z^2}+\alpha e^{-z^2}.
 \label{eq:v6-weinberg-counterfamily}
\end{equation}
Их нечётная часть одинакова для всех $\alpha$:
\begin{equation}
 g_\alpha^{\rm odd}(z)=ze^{-z^2},
\end{equation}
тогда как чётная часть равна
\begin{equation}
 g_\alpha^{\rm even}(z)=\alpha e^{-z^2}.
\end{equation}
Численный функциональный аудит дал один и тот же кинетический момент
$0.3133285343$, тогда как относительный коэффициент Вайнберга равен
$0,1,2$ при $\alpha=0,1,2$. Остаток совпадения нечётной части равен
$2.39\cdot10^{-17}$.

\begin{theorem}[Независимость коэффициента паринга]
Текущая кинетическая нормировка не определяет чётный нескалярный момент
фермионного спектрального действия. Поэтому коэффициент оператора
Вайнберга может непрерывно меняться и обращаться в ноль при неизменной
обычной фермионной кинетике.
\end{theorem}

\section{Почему вес одна седьмая не является массой}

Пусть один общий вес $w$ умножает оба члена:
\begin{equation}
 S=w\left[Z_\psi\overline\psi D\psi+
 \mu\psi^TC\psi\right].
\end{equation}
После нормировки поля физическая масса равна
\begin{equation}
 m=\frac{\mu}{Z_\psi}
\end{equation}
и не зависит от $w$. Аудит при $w=1/7,2/7,1$ дал одно значение
$0.01764705882$ с нулевым разбросом. Таким образом, выведенный вес
$1/7$ выбирает канал и его кратность, но не энергетический масштаб.

Литературный заряженно-лептонный множитель также сжимается только до
скаляра:
\begin{equation}
 P_0K_eP_0=\kappa_eP_0.
\end{equation}
Для двух допустимых контрольных $K_e$ получены разные значения
$\kappa_e=14/3$ и $2$. Направление одинаково, амплитуда различна.

\section{Окончательная граница}

Остаточная формула имеет структуру
\begin{equation}
 m_\nu\sim r_\tau\,\kappa_e\,\frac{v_H^2}{\Lambda}.
\end{equation}
Текущий родитель не фиксирует ни $r_\tau$, ни $\kappa_e$, ни
$\Lambda$, ни абсолютный $v_H$. Кроме того, однородный эффективный член
не выводит пространственный профиль смены ранга $0\to1$.

\begin{nogocriterion}[Запрет численной подстановки]
Нельзя использовать наблюдаемую нейтринную массу для выбора чётной части
спектральной функции, а затем объявлять результат предсказанием.
Контрсемейство \eqref{eq:v6-weinberg-counterfamily} доказывает, что такой
выбор является дополнительным параметром меры.
\end{nogocriterion}

\begin{proposition}[Частичное родительское происхождение]
Существующая архитектура фиксирует gauge-тип оператора, слабый
квадратичный носитель и одну голономно сжатую семейную ось. Она не
фиксирует полный семейный тензор, коэффициент, масштаб или локализацию.
Поэтому паринг Вайнберга допустим, но не выведен как параметрически
замкнутый член родителя.
\end{proposition}

\section{Следующий гейт}

Дальнейшее варьирование спектральной функции внутри текущей архитектуры
остановлено. Гейт
\texttt{version6\_spectral\_transition\_rank\_change\_localization\_gate}
должен вернуться к геометрически различающему вопросу: может ли уже
существующая динамика Хиггса и дефекта локализовать область, в которой
$W_\nu(H)$ меняет ранг $0\to1$, даже если абсолютная нейтринная масса
остаётся внешней.

\begin{protocol}
\begin{enumerate}
 \item gauge-тип оператора Вайнберга: \textbf{допустим};
 \item слабое направление: \textbf{фиксировано $B_\nu(H)$};
 \item размерность $C_3$-инвариантных семейных тензоров:
 \textbf{два};
 \item направление после $P_0$-сжатия: \textbf{одно};
 \item общая семейная амплитуда: \textbf{не фиксирована};
 \item коэффициент чётного спектрального момента:
 \textbf{не фиксирован};
 \item вес $1/7$ задаёт массу: \textbf{нет};
 \item параметрически чистая нейтринная масса: \textbf{не получена};
 \item локализованный переход ранга: \textbf{не выведен};
 \item физическое замыкание: \textbf{не достигнуто}.
\end{enumerate}
\end{protocol}

Машинный сертификат сохранён в
\path{s2t_v6_spectral_transition_weinberg_pairing_parent_gate.py} и
\path{s2t_v6_spectral_transition_weinberg_pairing_parent_gate_results.json}.

\begin{thebibliography}{9}
\bibitem{weinberg-parent-weinberg1979}
S.~Weinberg,
\emph{Baryon- and Lepton-Nonconserving Processes},
Physical Review Letters 43 (1979), 1566--1570.

\bibitem{weinberg-parent-sakellariadou-sitarz2019}
M.~Sakellariadou, A.~Sitarz,
\emph{Fermionic Spectral Action and the Origin of Nonzero Neutrino Masses},
Physics Letters B 795 (2019), 351--355;
arXiv:1903.09149.
\end{thebibliography}
